%% ch08.tex — Chapter 8: Data Pipelines, Data Lakes, and Real-Time Integration
%% Part of "Data Engineering" textbook

\chapter{Data Pipelines, Data Lakes, and Real-Time Integration}
\label{ch:pipelines}

\epigraph{\textit{``I believe that a unified log is the central data structure
that can solve the coherency problem in a distributed, heterogeneous data system
--- a structure so simple that we take it for granted, and so fundamental that
we underestimate it.''}}{Jay Kreps, \textit{The Log: What Every Software
Engineer Should Know About Real-Time Data's Unifying Abstraction},
LinkedIn Engineering, 2013}

\vspace{4pt}

\begin{chapterlearning}
After completing this chapter, you will be able to:
\begin{enumerate}
  \item Explain why traditional point-to-point ETL produces an $O(N^2)$
        integration complexity problem and describe how a central message
        bus reduces this to $O(N)$.
  \item Compare the five major data integration tools---Apache Sqoop,
        Apache Flume, Apache Kafka, Apache NiFi, and Spark Streaming---on
        their primary use case, data direction, and throughput characteristics.
  \item Describe how Apache Sqoop uses MapReduce to parallelise bulk data
        transfer between an RDBMS and HDFS, and configure an incremental
        import with \texttt{--incremental append}.
  \item Explain the three components of an Apache Flume agent---Source,
        Channel, and Sink---and distinguish between Memory Channel and
        File Channel in terms of durability guarantees.
  \item Define the five core Kafka concepts---Topic, Partition, Offset,
        Broker, and Consumer Group---and explain how they combine to deliver
        durability, parallelism, and exactly-once-capable messaging.
  \item Explain how Kafka achieves high throughput through sequential disk
        writes, zero-copy transfer via the Linux \texttt{sendfile()} system
        call, producer batching, and partition-level parallelism.
  \item Describe the three layers of the Lambda architecture---Batch Layer,
        Speed Layer, and Serving Layer---and identify its primary engineering
        disadvantage: the dual-code-path maintenance burden.
  \item Explain how the Kappa architecture eliminates the batch layer by
        using Kafka log replay for historical reprocessing, and identify
        workloads for which Kappa is \emph{not} a viable replacement for Lambda.
  \item Describe the four-layer architecture of a data lake---Ingestion,
        Storage, Processing, and Catalogue/Governance---and explain what
        turns a data lake into a data swamp.
  \item Explain how table formats such as Delta Lake and Apache Iceberg add
        ACID transaction semantics to object-store data lakes, and describe
        the role of the transaction log in enabling time-travel queries.
\end{enumerate}
\end{chapterlearning}

\bigskip

Chapter~\ref{ch:spark} introduced Apache Spark as an in-memory distributed
compute engine and Chapter~\ref{ch:sql-scale} extended Spark SQL into the
interactive analytics domain with Hive, Presto, and LLAP. Both chapters
addressed \emph{how to process} data that is already resident in HDFS or an
object store. This chapter asks the prior question: \emph{how does data arrive
there in the first place, and how can we ensure it arrives continuously, at
low latency, and with sufficient quality to be trusted downstream?}

The shift from nightly batch ETL to continuous data pipelines is one of the
most consequential architectural transitions in enterprise data engineering.
That transition was enabled by a deceptively simple idea: replace the web of
point-to-point connections between systems with a single, durable,
append-only log. Apache Kafka, introduced by LinkedIn engineers in 2011,
is the canonical implementation of that idea, and it has become the backbone
of modern event-driven architecture at companies ranging from Netflix and
Uber to bKash and Pathao.

The chapter covers the full pipeline picture: bulk batch transfer with
Sqoop and Flume; real-time streaming with Kafka; the two competing
system architectures---Lambda and Kappa---that organise batch and stream
processing into coherent pipelines; and the data lake as the storage
substrate that absorbs all of this data at petabyte scale. The chapter
closes with the governance problem that makes data lakes fail in practice,
and a preview of the table formats (Delta Lake, Apache Iceberg) that
are beginning to solve it.

%%====================================================================
\section{From Batch ETL to Streaming Pipelines}
\label{sec:etl-to-pipelines}
%%====================================================================

For two decades, the standard approach to moving data between enterprise
systems was the \term{ETL pipeline}: a scheduled job that \textbf{E}xtracts
data from a source system, \textbf{T}ransforms it into the target schema,
and \textbf{L}oads it into a destination (typically a data warehouse). ETL
pipelines are straightforward to build and reason about---a single job has
a defined source, a defined transformation, and a defined target---but they
scale poorly as the number of systems in an organisation grows.

%--------------------------------------------------------------------
\subsection{The $N^2$ Integration Problem}
\label{subsec:n2-problem}
%--------------------------------------------------------------------

In a point-to-point ETL architecture, each pair of systems that needs to
exchange data requires a dedicated ETL job. With $N$ systems, the number
of direct connections in the worst case is $N(N-1)/2 = O(N^2)$.

A bank with 50 operational systems (core banking, CRM, risk management,
fraud detection, mobile wallet, settlement, audit, compliance\ldots) may
require up to $50 \times 49 / 2 = 1{,}225$ separate ETL jobs. This
creates four interconnected problems:

\begin{enumerate}
  \item \textbf{Brittleness:} When the schema of source system $A$
        changes, every ETL job that reads from $A$ must be updated
        simultaneously. A column rename cascades into a multi-day
        incident.
  \item \textbf{No replay:} Traditional ETL jobs delete data from the
        source queue after consumption. If a downstream system is
        offline when a job runs, that data is lost unless the ETL
        job is re-run from scratch.
  \item \textbf{High latency:} All jobs are scheduled nightly (or
        hourly at best). Real-time fraud detection, live dashboards,
        and event-driven personalisation are architecturally impossible
        in this model.
  \item \textbf{Tight coupling:} Source and consumer are tightly bound.
        Adding a new consumer of an existing source requires modifying
        or duplicating the source-side ETL job.
\end{enumerate}

\begin{keyinsight}{The Event Bus Reduces $O(N^2)$ to $O(N)$}
Replacing direct ETL connections with a central \emph{event bus}
eliminates the coupling between producers and consumers. Each of the
$N$ systems connects to the bus once: producers write events to named
topics; consumers subscribe to the topics they care about. Adding a
new consumer requires \textbf{zero changes} to any producer. The total
number of connections is $2N$ (one publish connection and one subscribe
connection per system), reducing integration complexity from $O(N^2)$
to $O(N)$.
\end{keyinsight}

The logical architecture of a central event bus is shown in
Figure~\ref{fig:event-bus}.

\begin{figure}[htbp]
\centering
\begin{tikzpicture}[
    sys/.style={circle, draw=DEBlue!40, fill=DEBlue!12,
                minimum size=1.0cm, font=\small\bfseries\sffamily},
    bus/.style={rectangle, rounded corners=5pt, fill=DEAccent!18,
                draw=DEAccent!60, minimum width=2.6cm,
                minimum height=4.5cm, font=\small\bfseries\sffamily,
                align=center},
    arr/.style={-{Stealth[length=4pt,width=4pt]}, thick}
  ]
  %% Left panel: point-to-point
  \node[font=\small\bfseries\sffamily, text=DERed] at (-3.6, 5.0)
        {Point-to-Point ETL ($O(N^2)$)};
  \node[sys] (a1) at (-5.0, 4.0) {A};
  \node[sys] (a2) at (-2.2, 4.0) {B};
  \node[sys] (a3) at (-5.0, 2.4) {C};
  \node[sys] (a4) at (-2.2, 2.4) {D};
  \foreach \x in {a1,a2,a3,a4} {
    \foreach \y in {a1,a2,a3,a4} {
      \draw[DERed!30, very thin] (\x)--(\y);
    }
  }
  \node[font=\small\itshape, text=DERed!70] at (-3.6, 1.5)
        {fragile, no replay};

  %% Right panel: event bus
  \node[font=\small\bfseries\sffamily, text=DEGreen] at (3.6, 5.0)
        {Event Bus ($O(N)$)};
  \node[bus] (bus) at (3.6, 3.2) {Event\\Bus\\{\small (Kafka)}};

  \foreach \y/\lbl in {4.2/Src~A, 3.2/Src~B, 2.2/Src~C} {
    \node[sys, fill=DEGreen!20] (p\lbl) at (1.6, \y) {\scriptsize\lbl};
    \draw[arr, DEGreen!70] (p\lbl)--(bus);
  }
  \foreach \y/\lbl in {4.2/Con~1, 3.2/Con~2, 2.2/Con~3} {
    \node[sys, fill=DEBlue!20] (c\lbl) at (5.6, \y) {\scriptsize\lbl};
    \draw[arr, DEBlue!60] (bus)--(c\lbl);
  }
  \node[font=\small\itshape, text=DEGreen!70!black] at (3.6, 1.5)
        {decoupled, replayable};
\end{tikzpicture}
\caption{Point-to-point ETL (left) versus a central event bus (right).
         With $N=4$ systems, point-to-point requires up to six direct
         connections; the event bus requires only $2N = 8$ connections,
         and each source is oblivious to the number of consumers.}
\label{fig:event-bus}
\end{figure}

%--------------------------------------------------------------------
\subsection{The Data Integration Tool Landscape}
\label{subsec:tool-landscape}
%--------------------------------------------------------------------

The transition from batch ETL to streaming pipelines did not happen
overnight: a family of tools emerged to fill specific niches on the
spectrum between pure batch and pure stream processing.

\begin{table}[htbp]
\centering
\renewcommand{\arraystretch}{1.35}
\begin{tabular}{lllp{4.2cm}l}
  \toprule
  \textbf{Tool} & \textbf{Direction} & \textbf{Type}
    & \textbf{Primary Use Case} & \textbf{Throughput} \\
  \midrule
  Apache Sqoop
    & RDBMS $\leftrightarrow$ HDFS
    & Batch
    & Nightly bulk transfer of relational tables
    & GB--TB/hour \\
  Apache Flume
    & Logs $\to$ HDFS
    & Batch/micro-batch
    & Web server log ingestion; multi-source aggregation
    & MB--GB/hour \\
  Apache Kafka
    & Any $\to$ Any
    & Real-time stream
    & Centralised event bus; fraud, clickstream, IoT
    & GB--TB/hour \\
  Apache NiFi
    & Any $\to$ Any
    & Visual pipeline
    & No-code data routing; IoT edge to cloud
    & MB--GB/hour \\
  Spark Streaming
    & Kafka $\to$ HDFS/DB
    & Stream processing
    & Stateful aggregation; enrichment; stream joins
    & GB/hour \\
  \midrule
  \multicolumn{5}{l}{\small \textit{Rule of thumb:}
    Sqoop = bulk DB sync;\;
    Flume = log collection;\;
    Kafka = real-time events;\;
    NiFi = visual routing;\;
    Spark = stream computation} \\
  \bottomrule
\end{tabular}
\caption{Data integration tool landscape. No single tool is optimal for
         all scenarios; the correct choice depends on data direction, latency
         requirement, and operational complexity budget.}
\label{tab:tool-landscape}
\end{table}

%%====================================================================
\section{Batch Integration: Apache Sqoop and Apache Flume}
\label{sec:batch-integration}
%%====================================================================

Before real-time messaging systems like Kafka were available, two
Hadoop-ecosystem tools dominated data ingestion: \tool{Apache Sqoop}
for bulk database transfer and \tool{Apache Flume} for log collection.
Both remain in production at organisations with legacy Hadoop deployments,
and understanding their design illuminates why Kafka eventually displaced
them for most use cases.

%--------------------------------------------------------------------
\subsection{Apache Sqoop: Bulk RDBMS--HDFS Transfer}
\label{subsec:sqoop}
%--------------------------------------------------------------------

\term{Apache Sqoop} (\textit{SQL-to-Hadoop}) is a command-line tool that
transfers data in bulk between a relational database---MySQL, Oracle,
PostgreSQL, SQL Server, or any JDBC-compliant store---and HDFS (or Hive,
HBase). Its design insight is deceptively simple: convert a table scan into
a MapReduce job, where each Map task reads a non-overlapping range of rows
from the database in parallel using JDBC.

\textbf{Import pipeline.} When a Sqoop import command is submitted, Sqoop:
\begin{enumerate}
  \item Connects to the database via JDBC and queries the table metadata
        (column names, types, primary key range).
  \item Determines the number of mappers ($m$, default 4). It divides the
        primary key range $[k_{\min}, k_{\max}]$ into $m$ equal-width
        sub-ranges: mapper $i$ reads rows where
        $k \in [k_{\min} + (i-1)\cdot\delta,\; k_{\min} + i\cdot\delta)$.
  \item Submits a MapReduce job. Each Map task opens its own JDBC connection,
        executes the bounded \texttt{SELECT}, and writes its partition to
        HDFS as CSV, Avro, Parquet, or SequenceFile.
  \item There is no Reduce phase: each mapper writes an independent part
        file, and HDFS assembles them into the output directory.
\end{enumerate}

\begin{lstlisting}[style=shell, caption={Sqoop: full import and incremental import}]
# Full import: entire table to HDFS as Parquet, 8 parallel mappers
sqoop import \
  --connect jdbc:mysql://db-host:3306/finance \
  --username analyst \
  --password-file /user/analyst/.sqoop-pass \
  --table transactions \
  --as-parquetfile \
  --num-mappers 8 \
  --target-dir /data/raw/transactions \
  --delete-target-dir

# Incremental import: only rows newer than the last checkpoint
# Use this for nightly appends to an existing HDFS dataset
sqoop import \
  --connect jdbc:mysql://db-host:3306/finance \
  --username analyst \
  --password-file /user/analyst/.sqoop-pass \
  --table transactions \
  --as-parquetfile \
  --num-mappers 4 \
  --target-dir /data/raw/transactions \
  --incremental append \
  --check-column txn_timestamp \
  --last-value '2024-06-11 23:59:59'
\end{lstlisting}

\begin{defbox}{Sqoop Incremental Import}
An \textbf{incremental import} transfers only rows added since the previous
run, identified by a monotonically increasing column (\texttt{--check-column})
and a checkpoint value (\texttt{--last-value}). Two modes are supported:
\texttt{append} (new rows, identified by value greater than the checkpoint)
and \texttt{lastmodified} (rows whose modification timestamp exceeds the
checkpoint). Sqoop saves the new checkpoint to a \emph{saved job} so the
next run picks up where the previous left off.
\end{defbox}

\textbf{Export pipeline.} Sqoop can also export data \emph{from} HDFS
\emph{to} a relational database---for example, after a Spark batch job
produces a summary table on HDFS, Sqoop exports it back to a MySQL
reporting database for a business intelligence tool to query.

\begin{lstlisting}[style=shell, caption={Sqoop: export from HDFS to MySQL}]
sqoop export \
  --connect jdbc:mysql://db-host:3306/reports \
  --username loader \
  --password-file /user/loader/.sqoop-pass \
  --table daily_summary \
  --export-dir /data/processed/daily_summary \
  --input-fields-terminated-by '\t' \
  --update-mode allowinsert \
  --update-key summary_date,region
\end{lstlisting}

\begin{techdebt}{Why Sqoop Is Being Retired in Modern Stacks}
Sqoop's dependency on MapReduce makes it incompatible with YARN-free
deployments and cloud-native object stores where MapReduce locality
is irrelevant. Its JDBC-per-mapper approach hammers source databases
with concurrent connections, often requiring a dedicated read replica.
In modern data platforms, \tool{Kafka Connect} with a JDBC source
connector has replaced Sqoop for most use cases: it supports
streaming incremental ingestion (capturing changes via CDC), requires no
MapReduce infrastructure, and integrates with Kafka's replay semantics.
Sqoop reached End of Life in 2021 and is no longer receiving Apache
releases, though it remains in production at many Hadoop-era deployments.
\end{techdebt}

%--------------------------------------------------------------------
\subsection{Apache Flume: Log Ingestion Pipelines}
\label{subsec:flume}
%--------------------------------------------------------------------

\term{Apache Flume} is a distributed, reliable service for streaming log
data from multiple sources into HDFS. Its primary use case is collecting
semi-structured logs---web server access logs, application error logs,
IoT sensor readings---and storing them in HDFS for subsequent batch
analysis.

\textbf{The Flume agent model.} The fundamental unit of Flume is the
\term{agent}: a JVM process with three components.

\begin{defbox}{Flume Agent: Source, Channel, Sink}
A \textbf{Flume agent} consists of three components that form a directed
pipeline:
\begin{itemize}
  \item \textbf{Source}: receives events from an external system. Types
        include Exec source (\texttt{tail -F logfile}), Syslog source
        (UDP/TCP syslog), Avro source (from another Flume agent), and
        HTTP source (REST push).
  \item \textbf{Channel}: buffers events between Source and Sink.
        \emph{Memory Channel} (fast, non-durable) or \emph{File Channel}
        (writes to local disk; survives agent restart at the cost of
        throughput).
  \item \textbf{Sink}: writes buffered events to the destination.
        HDFS Sink (writes Avro or text files, rolling by time or size),
        Kafka Sink (forwards events to a Kafka topic), HBase Sink.
\end{itemize}
Events flow Source $\to$ Channel $\to$ Sink. The channel decouples
source read speed from sink write speed, absorbing bursts.
\end{defbox}

Multiple agents can be chained together: a \emph{fan-in} topology
aggregates logs from many web servers into a single \emph{collector agent}
before writing to HDFS, reducing the number of small files that would
result from each server writing directly to HDFS.

\begin{lstlisting}[style=shell, caption={Flume: agent configuration for web log collection}]
# flume-agent.conf
# Source: tail the nginx access log
agent.sources  = nginx_src
agent.channels = file_channel
agent.sinks    = hdfs_sink

agent.sources.nginx_src.type           = exec
agent.sources.nginx_src.command        = tail -F /var/log/nginx/access.log
agent.sources.nginx_src.channels       = file_channel

agent.channels.file_channel.type       = file
agent.channels.file_channel.checkpointDir = /var/flume/checkpoint
agent.channels.file_channel.dataDirs      = /var/flume/data

agent.sinks.hdfs_sink.type             = hdfs
agent.sinks.hdfs_sink.hdfs.path        = hdfs://namenode:9000/logs/%Y/%m/%d
agent.sinks.hdfs_sink.hdfs.fileType    = DataStream
agent.sinks.hdfs_sink.hdfs.rollInterval= 300
agent.sinks.hdfs_sink.hdfs.rollSize    = 134217728
agent.sinks.hdfs_sink.channel          = file_channel
\end{lstlisting}

\begin{caution}{Small Files and the HDFS Namespace Tax}
Flume's rolling policy controls how often a new HDFS file is created.
A short roll interval (e.g., every 10 seconds) produces thousands of
small files per day, each consuming a NameNode metadata entry and
degrading NameNode heap. Best practice is to use a roll interval of
300--600 seconds and a roll size of 128~MB (one HDFS block), so each
rolled file is exactly one block. Alternatively, a downstream compaction
job (Spark or Hive \texttt{INSERT OVERWRITE \ldots\ SELECT \ldots}) merges
small files nightly into large Parquet files.
\end{caution}

\begin{systemdesign}{Flume Fan-In: Aggregating Logs from Many Servers}
A typical Flume deployment for a web application with 20 application
servers uses two tiers: twenty \emph{edge agents} (one per server)
each with an Exec or Syslog source and a Memory Channel writing to
an Avro Sink; and two \emph{collector agents} each with an Avro Source
(receiving from 10 edge agents) and a File Channel writing to an HDFS
Sink. The collector tier provides durability (File Channel survives
collector restart) and aggregation (reducing 20 HDFS writers to 2,
greatly reducing small file pressure). This fan-in topology is specified
entirely in the Flume agent configuration files---no code is required.
\end{systemdesign}

%%====================================================================
\section{Apache Kafka: The Distributed Append-Only Log}
\label{sec:kafka}
%%====================================================================

Apache Kafka was created by Jay Kreps, Neha Narkhede, and Jun Rao at
LinkedIn in 2010--2011 to solve the $N^2$ integration problem that LinkedIn
faced with its rapidly growing activity data pipeline. It was open-sourced
to the Apache Software Foundation in 2012 and has since become the dominant
event streaming platform in the industry.

Kafka's central innovation is not a clever algorithm or a novel hardware
optimisation: it is a \emph{rethinking of what a message broker should be}.
Traditional message queues (ActiveMQ, RabbitMQ) model messages as transient
objects: once a consumer acknowledges receipt, the message is deleted.
Kafka models messages as \emph{records in a durable, immutable, ordered
log}: they are retained for a configurable period regardless of whether
they have been consumed, and any consumer can read them at any time by
specifying an offset.

%--------------------------------------------------------------------
\subsection{The Core Log Abstraction}
\label{subsec:kafka-log}
%--------------------------------------------------------------------

A \term{log}, in the computer science sense, is the simplest possible
storage abstraction: an append-only, totally ordered sequence of records.
Every database transaction log, every write-ahead log, every file system
journal is a log in this sense. Kafka's insight is that this same
abstraction---which databases use for durability---can serve as the
backbone of an entire distributed data infrastructure.

\begin{defbox}{Kafka Log Abstraction}
In Kafka, a \textbf{log} is a durable, append-only sequence of records.
Each record has:
\begin{itemize}
  \item A \textbf{key} (optional): a byte array used to determine partition
        assignment (all records with the same key go to the same partition,
        preserving key-level ordering).
  \item A \textbf{value}: the payload bytes (serialised JSON, Avro, Protobuf,
        or raw bytes).
  \item A \textbf{timestamp}: producer-assigned or broker-assigned event time.
  \item An \textbf{offset}: a monotonically increasing integer assigned by the
        broker that uniquely identifies the record within its partition.
\end{itemize}
Records are \textbf{never modified or deleted} within the retention period.
The log is the single source of truth; all consumers derive their view
from it by tracking their own read offset.
\end{defbox}

Figure~\ref{fig:kafka-log} contrasts the traditional queue model with
Kafka's log model. In a traditional queue, a message is removed after
it is consumed; a late-joining consumer has no access to past messages.
In Kafka, all messages are retained for a configurable period, enabling
any consumer---including one that joins weeks after the messages were
written---to read from offset 0 and replay the full history.

\begin{figure}[htbp]
\centering
\begin{tikzpicture}[
    rec/.style={rectangle, draw=DEBlue!35, fill=DEBlue!10,
                minimum width=0.8cm, minimum height=0.65cm,
                font=\scriptsize\sffamily},
    arr/.style={-{Stealth[length=4pt,width=4pt]}, thick}
  ]
  %% Traditional queue
  \node[font=\small\bfseries\sffamily, text=DERed] at (0, 4.8)
        {Traditional Queue};
  \foreach \i/\x in {0/-1.5, 1/-0.5, 2/0.5, 3/1.5} {
    \node[rec, fill=DERed!15, draw=DERed!30] at (\x, 4.1) {\tiny msg};
  }
  \node[font=\scriptsize, text=DERed!70, align=center] at (0, 3.5)
        {consumed = deleted\\late consumer loses data};

  %% Kafka log
  \node[font=\small\bfseries\sffamily, text=DEBlue] at (7.0, 4.8)
        {Kafka Partition (append-only log)};
  \foreach \off/\x in {0/4.5, 1/5.4, 2/6.3, 3/7.2, 4/8.1, 5/9.0} {
    \node[rec] (k\off) at (\x, 4.1) {\tiny\off};
  }
  \draw[arr, DEGray!50] (4.1,4.1)--(4.1,4.5)
        node[above, font=\scriptsize\sffamily]{oldest};
  \draw[arr, DEAccent] (9.4,4.1)--(9.4,4.5)
        node[above, font=\scriptsize\sffamily]{newest};

  %% Consumer A at offset 2
  \draw[arr, DEGreen!70, thick] (6.3, 3.5)--(6.3, 3.8)
        node[below=4pt, font=\scriptsize\sffamily\itshape,
             text=DEGreen!70!black, align=center]{Con A\\off=2};
  %% Consumer B at offset 5
  \draw[arr, DEBlue!60, thick] (9.0, 3.5)--(9.0, 3.8)
        node[below=4pt, font=\scriptsize\sffamily\itshape,
             text=DEBlue, align=center]{Con B\\off=5};

  \node[font=\scriptsize\itshape, text=DEGray, align=center] at (7.0, 2.9)
        {each consumer tracks its own offset\\
         messages retained for days or weeks};
\end{tikzpicture}
\caption{Traditional queue versus Kafka's append-only log. In the queue model,
         consumed messages are deleted. In Kafka, all messages are retained;
         each consumer independently tracks its read offset and can replay
         history by resetting to any past offset.}
\label{fig:kafka-log}
\end{figure}

%--------------------------------------------------------------------
\subsection{Kafka Architecture: Topics, Partitions, Offsets, and Brokers}
\label{subsec:kafka-arch}
%--------------------------------------------------------------------

Kafka's architecture has five foundational concepts that work together to
deliver durability, parallelism, and horizontal scalability.

\begin{enumerate}
  \item \textbf{Topic.} A \term{topic} is a named, durable, append-only log
        for a category of events. Topics are analogous to database tables
        but append-only: \texttt{bkash-transactions}, \texttt{pathao-gps-pings},
        and \texttt{dhaka-weather-sensors} might each be a Kafka topic.
        Producers write to a topic; consumers read from it.

  \item \textbf{Partition.} Each topic is split into $P \geq 1$ ordered
        \term{partitions}. A partition is an independent ordered log residing
        on a single broker. Partitions provide two things:
        \emph{parallelism} (multiple producers can write to different
        partitions simultaneously; multiple consumers in a group can
        each process a distinct partition) and
        \emph{key-based ordering} (all messages with the same key are
        routed to the same partition via key hashing, guaranteeing that
        a given customer's transactions are always read in order).

  \item \textbf{Offset.} Each message within a partition has a unique,
        monotonically increasing integer \term{offset}---its position in
        the partition log. Offsets within a partition are dense and
        ordered; they are not globally unique across partitions. Each
        consumer independently tracks its read offset; the broker does
        not maintain per-consumer state (offsets are stored by consumers
        in a special Kafka topic named \texttt{\_\_consumer\_offsets}).

  \item \textbf{Broker.} A \term{broker} is a Kafka server process that
        stores partitions on its local disk and serves producers and
        consumers over the Kafka wire protocol. A Kafka cluster contains
        $B$ brokers. Each partition has one \term{leader replica} (handles
        all reads and writes) and $R-1$ \term{follower replicas} (copy
        data from the leader asynchronously; stand in as leader on
        failure). The default replication factor is $R=3$, tolerating
        the simultaneous failure of any two brokers without data loss.

  \item \textbf{ZooKeeper / KRaft.} Kafka versions prior to 3.0 used
        Apache ZooKeeper for cluster metadata and leader election.
        Kafka 3.0 introduced \tool{KRaft} (Kafka Raft), a self-managed
        metadata quorum built directly into the Kafka brokers, removing
        the ZooKeeper dependency and simplifying deployment.
\end{enumerate}

\begin{figure}[htbp]
\centering
\begin{tikzpicture}[
    part/.style={rectangle, draw=DEBlue!35, fill=DEBlue!10,
                 minimum width=2.8cm, minimum height=0.6cm,
                 font=\small\sffamily, rounded corners=2pt},
    brk/.style={rectangle, rounded corners=5pt, draw=DEAccent!60,
                fill=DEAccent!10, minimum width=3.0cm,
                minimum height=2.6cm},
    arr/.style={-{Stealth[length=4pt,width=4pt]}, thick}
  ]
  %% Topic label
  \node[font=\normalsize\bfseries\sffamily, text=DEBlue] at (3.6, 6.4)
        {Topic: \texttt{bkash-txns}\ \ (3 partitions, RF\,=\,2)};

  %% Brokers
  \node[brk] at (0.8, 4.6) {};
  \node[brk] at (6.4, 4.6) {};
  \node[font=\small\bfseries\sffamily, text=DEAccent!80!black] at (0.8, 5.9)
        {Broker 1};
  \node[font=\small\bfseries\sffamily, text=DEAccent!80!black] at (6.4, 5.9)
        {Broker 2};

  %% Partitions on Broker 1
  \node[part, fill=DEBlue!25, text=white] (p0l) at (0.8, 5.2)
        {P0 (leader)};
  \node[part]                             (p1l) at (0.8, 4.5)
        {P1 (leader)};
  \node[part]                             (p2r) at (0.8, 3.8)
        {P2 (replica)};

  %% Partitions on Broker 2
  \node[part]                             (p0r) at (6.4, 5.2)
        {P0 (replica)};
  \node[part, fill=DEBlue!25, text=white] (p2l) at (6.4, 4.5)
        {P2 (leader)};
  \node[part]                             (p1r) at (6.4, 3.8)
        {P1 (replica)};

  %% Replication arrows
  \draw[arr, DERed!40, dashed, bend left=15] (p0l) to (p0r);
  \draw[arr, DERed!40, dashed, bend right=15] (p1l) to (p1r);
  \draw[arr, DERed!40, dashed, bend left=10] (p2l) to (p2r);
  \node[font=\scriptsize\itshape, text=DERed!60] at (3.6, 5.7)
        {replication};

  %% Producers
  \node[font=\small\bfseries\sffamily] at (-2.2, 5.2) {Producer};
  \draw[arr, DEGreen!70]  (-1.4,5.2)--(0.0,5.2)
        node[midway,above,font=\scriptsize\sffamily]{key=txn\_id};
  \draw[arr, DEGreen!50]  (-1.4,4.5)--(0.0,4.5);

  %% Consumer Group
  \node[font=\small\bfseries\sffamily] at (9.4, 5.2) {Consumer\\Group};
  \draw[arr, DEBlue!60] (7.2,5.2)--(8.6,5.2)
        node[midway,above,font=\scriptsize\sffamily]{analytics};
  \draw[arr, DEBlue!60] (7.2,4.5)--(8.6,4.5);
  \draw[arr, DEBlue!40] (7.2,3.8)--(8.6,3.8);
\end{tikzpicture}
\caption{Kafka cluster with two brokers, one topic, three partitions, and
         replication factor 2. Partition leaders handle all reads and writes;
         followers maintain copies for fault tolerance. If Broker~1 fails,
         Broker~2 promotes its P0 and P1 replicas to leaders.}
\label{fig:kafka-arch}
\end{figure}

%--------------------------------------------------------------------
\subsection{Producers, Consumers, and Consumer Groups}
\label{subsec:kafka-producers-consumers}
%--------------------------------------------------------------------

A \term{Kafka producer} is any application that publishes messages to a
topic. The producer selects the target partition either by hashing the
message key (\texttt{hash(key) \% P}) or by round-robin if no key is
specified. The producer batches messages in memory and sends a batch
to the broker's partition leader in a single network request, trading
latency for throughput via the \texttt{batch.size} and \texttt{linger.ms}
configuration parameters.

A \term{Kafka consumer} is any application that reads messages from one
or more topics by polling the broker and advancing its read offset. The
consumer never \emph{pushes} messages back to acknowledge deletion; it
simply records that it has processed up to a given offset.

\term{Consumer groups} are the key mechanism for parallel consumption.
A consumer group is a named set of consumer processes that collectively
consume all partitions of a topic. Kafka's partition-assignment guarantee
ensures that each partition is assigned to \emph{exactly one consumer}
in the group at any point in time. This means:
\begin{itemize}
  \item Adding consumers to a group (up to $P$ consumers) increases
        parallelism proportionally.
  \item Two independent applications (e.g., a fraud-detection service and
        an analytics service) reading the same topic use \emph{different}
        consumer groups. Each group maintains its own set of offsets,
        so both applications receive every message independently.
  \item When a consumer joins or leaves the group, Kafka triggers a
        \term{rebalance}: partitions are reassigned among the remaining
        consumers. Modern Kafka (2.4+) uses cooperative incremental
        rebalancing to minimise the disruption.
\end{itemize}

\begin{lstlisting}[style=pyspark, caption={Kafka producer and consumer in Python (kafka-python)}]
import json, random, time
from kafka import KafkaProducer, KafkaConsumer

BROKER = "localhost:9092"
TOPIC  = "bkash-transactions"

# ---------- Producer ----------
producer = KafkaProducer(
    bootstrap_servers=BROKER,
    key_serializer=lambda k: k.encode(),
    value_serializer=lambda v: json.dumps(v).encode(),
    # batch.size=16384 (16 KB), linger.ms=5 — trade latency for throughput
    batch_size=16384,
    linger_ms=5,
    acks='all',          # wait for all in-sync replicas
    retries=5,
)

for i in range(10_000):
    txn = {
        "txn_id":    f"TXN{i:08d}",
        "sender":    random.randint(1, 50_000),
        "receiver":  random.randint(1, 50_000),
        "amount":    round(random.uniform(10, 5000), 2),
        "currency":  "BDT",
        "timestamp": int(time.time() * 1000),
    }
    # key = str(sender): ensures all transactions from the same
    # sender land in the same partition, preserving send-order
    producer.send(TOPIC, key=str(txn["sender"]), value=txn)

producer.flush()   # block until all buffered records are sent
print("Produced 10,000 transactions")

# ---------- Consumer (fraud-detection group) ----------
consumer = KafkaConsumer(
    TOPIC,
    bootstrap_servers=BROKER,
    group_id="fraud-detection-v1",
    auto_offset_reset="earliest",         # start from offset 0 if new group
    enable_auto_commit=True,
    auto_commit_interval_ms=5000,
    value_deserializer=lambda m: json.loads(m.decode()),
    key_deserializer=lambda k: k.decode() if k else None,
)

for record in consumer:
    txn = record.value
    # Simple heuristic: flag large transactions for review
    if txn["amount"] > 4500:
        print(f"[ALERT] Large txn {txn['txn_id']}: "
              f"BDT {txn['amount']:.2f} "
              f"(partition={record.partition}, offset={record.offset})")
\end{lstlisting}

%--------------------------------------------------------------------
\subsection{Kafka Throughput: Sequential I/O and Zero-Copy Transfer}
\label{subsec:kafka-throughput}
%--------------------------------------------------------------------

The Kreps et al.\ (2011) paper reports Kafka producer throughput of
505{,}000 messages/sec---roughly $11\times$ that of ActiveMQ and $3.4\times$
that of RabbitMQ on equivalent hardware. Four engineering decisions explain
this performance advantage.

\begin{enumerate}
  \item \textbf{Sequential disk writes.} Kafka never updates messages in
        place. Partition logs are written by appending to the active
        segment file. Sequential writes on a spinning hard disk achieve
        500--600~MB/s; random writes to the same disk achieve only
        50--100~MB/s due to seek latency. Kafka exploits this difference
        fully. The OS page cache prefetches the next disk blocks
        automatically, making Kafka's effective write speed close to
        raw disk bandwidth even without SSDs.

  \item \textbf{Zero-copy transfer.} When a consumer fetches messages,
        a traditional broker copies data along the following path:
        disk $\to$ kernel buffer $\to$ user-space buffer $\to$ socket
        buffer $\to$ NIC (four copies, two context switches). Kafka
        uses the Linux \texttt{sendfile()} system call to transfer data
        from the kernel buffer directly to the NIC, bypassing the
        user-space copy entirely (two copies, zero context switches).
        This halves CPU usage for read-heavy workloads.

  \item \textbf{Batch compression.} Producers compress a \emph{batch}
        of messages as a single unit (Snappy, LZ4, Zstandard, or GZIP)
        before transmitting to the broker. The broker stores the
        compressed batch without decompressing it; consumers decompress
        on receipt. For structured log data, batch compression
        ratios of 5--10$\times$ are typical, reducing both network
        and disk I/O proportionally.

  \item \textbf{Partition parallelism.} Throughput scales linearly
        with partition count: $P$ partitions allow $P$ parallel
        producer threads and up to $P$ consumers in a group to
        operate simultaneously. Increasing partitions from 1 to 32
        increases end-to-end throughput by a factor approaching 32,
        bounded by broker and network capacity.
\end{enumerate}

\begin{keyinsight}{Kafka's Disk Is Faster Than Memory for Sequential I/O}
A common assumption is that memory-based messaging (e.g., keeping all
messages in RAM) is faster than disk-based messaging. This is true
for \emph{random} access. For the \emph{sequential} access pattern
that Kafka's append-only log produces, the OS page cache makes
sequential disk access \emph{faster than a heap-allocated data structure}
at large scale: sequential disk reads saturate the memory bus rather
than L3 cache, and the page cache is managed by the kernel's LRU policy
which is better tuned than a custom in-process cache for this workload.
\end{keyinsight}

%--------------------------------------------------------------------
\subsection{Retention, Replay, and Exactly-Once Semantics}
\label{subsec:kafka-retention}
%--------------------------------------------------------------------

\textbf{Retention policy.} Kafka retains messages for a configurable
period (\texttt{retention.ms}, default 7 days) or up to a configurable
total log size (\texttt{retention.bytes}). When both conditions are
reached, the oldest segment files are deleted. A consumer that has not
read a partition before its messages expire will miss those messages---a
design choice that prioritises throughput over guaranteed delivery to
slow consumers.

\textbf{Replay.} Because offsets are persistent and messages are not
deleted on consumption, any consumer can reset its offset to any point
in the retention window and re-read messages. This enables:
\begin{itemize}
  \item \textbf{Backfilling}: a new downstream service can be launched
        and immediately process historical data without re-ingesting
        from the source.
  \item \textbf{Bug recovery}: if a consumer has a bug that corrupts its
        state, it can reset to the last known-good offset and reprocess.
  \item \textbf{Multi-consumer independence}: a second consumer group
        reading the same topic starts fresh at offset 0 without affecting
        the first group's offset position.
\end{itemize}

\textbf{Exactly-once semantics.} Kafka 0.11 (2017) introduced two
mechanisms for exactly-once message delivery:
\begin{itemize}
  \item \textbf{Idempotent producer.} The producer assigns a monotonically
        increasing sequence number to each batch. If the broker receives
        a duplicate batch (e.g., due to a network retry), it deduplicates
        using the sequence number. This prevents duplicate writes caused
        by producer retries.
  \item \textbf{Transactional API.} A producer can open a transaction,
        send messages to multiple partitions, and either commit (making
        all messages visible atomically) or abort (rolling back). Combined
        with \texttt{isolation.level=read\_committed} on the consumer side,
        this provides exactly-once semantics across the full
        produce--consume--produce pipeline.
\end{itemize}

\begin{techdebt}{Schema Governance at Scale: The Schema Registry}
At small scale, Kafka consumers can deserialise message bytes using
any format they agree upon with producers. At scale---4{,}000 topics
and 400 microservices, as at Netflix---schema evolution breaks consumers:
a producer that adds a field to its Avro schema without warning causes
all consumers to crash on the next schema change.

The \tool{Confluent Schema Registry} (open-source) solves this by
registering every Avro, Protobuf, or JSON Schema version under a
subject (typically \texttt{\{topic\}-value}). Producers embed the
schema ID (an integer) in each message; consumers look up the schema
by ID on receipt. The Registry enforces a \emph{compatibility policy}
(BACKWARD, FORWARD, FULL) that rejects schema changes which would break
existing consumers. At Netflix and LinkedIn, deploying a schema registry
is considered non-negotiable for production Kafka usage.
\end{techdebt}

%%====================================================================
\section{Lambda Architecture}
\label{sec:lambda}
%%====================================================================

As real-time stream processing systems matured alongside batch processing
frameworks, practitioners faced a design dilemma: how to serve queries that
require both low latency (seconds) \emph{and} high accuracy (based on
complete historical data), when stream processors are fast but approximate
and batch processors are accurate but slow?

Nathan Marz formulated the \term{Lambda architecture} in 2011 as a systematic
answer to this dilemma. The name is drawn from the Greek letter $\lambda$,
reflecting the three-way branching of the data path.

%--------------------------------------------------------------------
\subsection{The Three-Layer Design}
\label{subsec:lambda-layers}
%--------------------------------------------------------------------

\begin{defbox}{Lambda Architecture}
The \textbf{Lambda architecture} is a data pipeline pattern that serves
both accurate historical batch views and low-latency real-time views from
the same input data, without choosing one at the expense of the other.
It has three layers:
\begin{enumerate}
  \item \textbf{Batch Layer}: stores the complete, immutable master dataset
        (typically on HDFS or S3). Periodically recomputes \emph{batch views}
        by running MapReduce, Spark, or Hive jobs over the entire dataset.
        Output is \emph{accurate} but \emph{high-latency}: views are
        1~hour to 24~hours stale.
  \item \textbf{Speed Layer}: processes new data arriving since the last
        batch run using a stream processor (Kafka + Spark Streaming or
        Flink). Produces \emph{real-time views} covering the most recent
        window. Output is \emph{low-latency} ($<$1~second to seconds) but
        may be \emph{approximate} (bounded by available state in the
        stream processor).
  \item \textbf{Serving Layer}: merges batch views and real-time views
        to answer user queries. A query result equals
        \emph{batch view (up to N hours ago)} $\cup$ \emph{real-time view
        (last N hours)}. Apache Druid, Cassandra, or a custom merge index
        typically implements the serving layer.
\end{enumerate}
\end{defbox}

\begin{figure}[htbp]
\centering
\begin{tikzpicture}[
    box/.style={rectangle, rounded corners=5pt, minimum width=3.6cm,
                minimum height=0.8cm, font=\small\bfseries\sffamily,
                align=center},
    arr/.style={-{Stealth[length=5pt,width=5pt]}, thick}
  ]
  %% Raw data
  \node[box, fill=DEBlue!80, text=white, minimum width=6.0cm]
        (raw) at (0, 6.0) {Raw Data (immutable event log)};

  %% Batch layer
  \node[box, fill=DERed!18, minimum width=2.6cm]
        (bl) at (-2.5, 4.6) {Batch Layer\\{\small HDFS + Spark/Hive}};
  %% Speed layer
  \node[box, fill=DEGreen!22, minimum width=2.6cm]
        (sl) at (2.5, 4.6) {Speed Layer\\{\small Kafka + Flink}};

  \draw[arr, DERed!55] (raw)--(-2.5,5.4)--(-2.5,5.1);
  \draw[arr, DEGreen!55] (raw)--(2.5,5.4)--(2.5,5.1);

  %% Views
  \node[box, fill=DERed!10, minimum width=2.6cm] (bv) at (-2.5, 3.3)
        {Batch Views\\{\small hours-old, accurate}};
  \node[box, fill=DEGreen!12, minimum width=2.6cm] (rv) at (2.5, 3.3)
        {Real-time Views\\{\small seconds-old, approx.}};

  \draw[arr, DERed!45]   (bl)--(bv);
  \draw[arr, DEGreen!45] (sl)--(rv);

  %% Serving layer
  \node[box, fill=DEAccent!22, minimum width=6.0cm]
        (srv) at (0, 2.0) {Serving Layer (merge + query)};
  \draw[arr, DERed!45]   (bv)--(srv);
  \draw[arr, DEGreen!45] (rv)--(srv);

  %% Application
  \node[box, fill=DEBlue!12, minimum width=6.0cm]
        (app) at (0, 0.8) {Application / Dashboard};
  \draw[arr, DEBlue!50] (srv)--(app);
\end{tikzpicture}
\caption{Lambda architecture. Raw data flows into both the batch layer
         (for accurate periodic recomputation) and the speed layer (for
         low-latency real-time views). The serving layer merges both view
         types to answer queries.}
\label{fig:lambda-arch}
\end{figure}

%--------------------------------------------------------------------
\subsection{Strengths and the Dual-Code-Path Problem}
\label{subsec:lambda-weaknesses}
%--------------------------------------------------------------------

The Lambda architecture has four significant strengths:
\begin{itemize}
  \item \textbf{Accuracy from immutability.} The batch layer always
        recomputes from the raw, immutable dataset. A bug in a past
        batch job can be corrected by fixing the logic and re-running
        the job---the historical truth is never lost.
  \item \textbf{Low latency for recent data.} The speed layer delivers
        query results for the most recent window within seconds,
        regardless of how large the historical dataset is.
  \item \textbf{Fault tolerance.} No derived view is irreplaceable:
        batch views and real-time views can both be recomputed from
        the master dataset or from Kafka's retained log.
  \item \textbf{Separation of concerns.} Batch and speed components
        can be built with the most appropriate tools for each: Hive
        for batch SQL, Flink for streaming CEP, for example.
\end{itemize}

The Lambda architecture's critical weakness is the \textbf{dual-code-path
problem}: the same business logic must be implemented \emph{twice}---once
in the batch layer (SQL or Spark batch) and once in the speed layer
(streaming SQL or Flink operators).

\begin{caution}{The Dual-Code-Path Maintenance Burden}
Consider a data engineering team computing ``daily active users per region''
for a fraud monitoring system. In a Lambda deployment this requires:
\begin{enumerate}
  \item A Spark batch job (running nightly) that aggregates the full
        historical table and writes batch views for periods up to midnight.
  \item A Flink streaming job that maintains a rolling count of active
        users from midnight to now (the real-time view).
\end{enumerate}
When the definition of ``active user'' changes (e.g., the business adds
a minimum session-duration threshold), \emph{both jobs must be updated
simultaneously}. Industry reports indicate that 30--40\% of data
engineering effort in Lambda deployments is spent keeping the two
implementations in sync---a tax that grows with the number of metrics
and the frequency of business logic changes.
\end{caution}

%%====================================================================
\section{Kappa Architecture}
\label{sec:kappa}
%%====================================================================

In 2014, Jay Kreps (Kafka co-creator) proposed the \term{Kappa
architecture} as a simplification of Lambda. The core argument: if a
stream processing system is powerful enough to reprocess historical data
at the same rate as live data, the batch layer is unnecessary.

%--------------------------------------------------------------------
\subsection{Stream-Only Processing and Log Replay}
\label{subsec:kappa-design}
%--------------------------------------------------------------------

\begin{defbox}{Kappa Architecture}
The \textbf{Kappa architecture} replaces Lambda's three-layer design
with two components:
\begin{enumerate}
  \item \textbf{Log storage} (Kafka with long-term retention, or an
        object store like S3 with Apache Iceberg providing table-level
        access): stores the complete immutable event history. This
        replaces the HDFS batch layer.
  \item \textbf{Stream processing engine} (Apache Flink or Spark
        Structured Streaming): the \emph{single code path} that handles
        both live events and historical replay. To reprocess history,
        a new consumer reads from Kafka offset 0 and processes events
        in order until it catches up to the live stream.
\end{enumerate}
The \textbf{reprocessing workflow} when business logic changes:
\begin{enumerate}
  \item Deploy a new version of the stream job alongside the current
        version (they write to different output topics or tables).
  \item Start the new job at Kafka offset~0; it reprocesses the full
        retained history.
  \item When the new job catches up to the live stream, atomically
        swap the serving layer to read from the new output.
  \item Shut down the old version.
\end{enumerate}
\end{defbox}

\begin{figure}[htbp]
\centering
\begin{tikzpicture}[
    box/.style={rectangle, rounded corners=5pt, minimum width=4.2cm,
                minimum height=0.8cm, font=\small\bfseries\sffamily,
                align=center},
    arr/.style={-{Stealth[length=5pt,width=5pt]}, thick}
  ]
  %% Raw events
  \node[box, fill=DEBlue!80, text=white]
        (raw) at (0, 5.8) {Raw Events (producers)};

  %% Kafka log
  \node[box, fill=DEAccent!22, minimum height=1.1cm]
        (kaf) at (0, 4.4) {Kafka (durable log)\\{\small history retained at offset 0\ldots N}};
  \draw[arr, DEBlue!55] (raw)--(kaf);

  %% Stream job live
  \node[box, fill=DEGreen!25]
        (live) at (0, 3.0) {Stream Job v2 (live, offset = now)};
  %% Stream job replay
  \node[box, fill=DEGreen!12]
        (rply) at (0, 2.0) {Stream Job v2 (replay, offset = 0)};

  \draw[arr, DEAccent!70] (kaf)--(live);
  \draw[arr, DEAccent!40, dashed] (kaf)--(rply);

  %% Serving store
  \node[box, fill=DEBlue!15]
        (srv) at (0, 0.8) {Serving Store (swap when caught up)};
  \draw[arr, DEGreen!60] (live)--(srv);
  \draw[arr, DEGreen!30, dashed] (rply)--
        node[right,font=\scriptsize\sffamily\itshape]{when ready} (srv);

  \node[font=\small\bfseries\sffamily, text=DEAccent!80!black] at (0, 0.2)
        {one code path — infinite replayability};
\end{tikzpicture}
\caption{Kappa architecture. A single stream processing job handles
         both live events and historical replay. When business logic
         changes, a new job version is started at offset~0, reprocesses
         history, and atomically replaces the old serving store.}
\label{fig:kappa-arch}
\end{figure}

%--------------------------------------------------------------------
\subsection{Lambda vs.\ Kappa: Comparison and Selection}
\label{subsec:lambda-vs-kappa}
%--------------------------------------------------------------------

\begin{table}[htbp]
\centering
\renewcommand{\arraystretch}{1.4}
\begin{tabular}{lp{4.8cm}p{4.8cm}}
  \toprule
  \textbf{Property}
    & \textbf{Lambda}
    & \textbf{Kappa} \\
  \midrule
  Code paths
    & Two (batch + stream), must stay in sync
    & One (stream only) \\
  Historical reprocessing
    & Re-run batch job on HDFS
    & Replay Kafka log from offset~0 \\
  Latency
    & Low (speed layer) $+$ High (batch layer)
    & Uniformly low \\
  Complexity
    & High: two systems, merge serving layer
    & Lower: single pipeline \\
  Batch throughput
    & Very high (Spark/Hive, full dataset in parallel)
    & Bounded by stream engine throughput \\
  Long-term storage cost
    & HDFS (commodity storage, dense packing)
    & Kafka retention (more expensive per GB than HDFS) \\
  Fault tolerance
    & Disk-based; individual task retries
    & Checkpoint-based; replay on failure \\
  Best when
    & Business logic is too complex for streaming SQL;
      regulatory batch reconciliation required;
      training datasets for ML must be reproducible
    & Stream engine is fast enough to reprocess history;
      single code path is worth the stream-storage cost;
      schema/logic evolves frequently \\
  \bottomrule
\end{tabular}
\caption{Lambda versus Kappa architecture. Neither is universally superior;
         the correct choice depends on the organisation's latency requirements,
         operational capabilities, and the pace of business logic change.}
\label{tab:lambda-vs-kappa}
\end{table}

\begin{caution}{When Kappa Is Not Suitable}
Kappa is \emph{not} a universal Lambda replacement. Three scenarios
favour Lambda:
\begin{enumerate}
  \item \textbf{Terabyte-scale reprocessing is too slow for the stream
        engine.} A Flink job reprocessing 5~PB of Kafka history at
        1~GB/s takes \textasciitilde{}58~days. A Spark batch job over
        the same data on a 200-node cluster completes in hours.
  \item \textbf{Batch SQL analytics dominate the workload.} If the
        primary consumers are data analysts running ad-hoc HiveQL or
        Spark SQL queries over months of history, a Kafka-backed Kappa
        architecture cannot serve these workloads without first landing
        the data to a queryable store---effectively recreating the
        batch layer.
  \item \textbf{Regulatory requirement mandates independent batch
        reconciliation.} Bangladesh Bank Circular~8 (2022) requires
        that all mobile financial service transactions be reconciled in
        a separate batch system to detect discrepancies between the
        real-time ledger and the daily settlement---a process that cannot
        be replaced by a streaming job.
\end{enumerate}
\end{caution}

%%====================================================================
\section{Data Lakes and the Data Swamp Risk}
\label{sec:data-lakes}
%%====================================================================

Both Lambda and Kappa architectures produce large volumes of raw and
processed data that must be stored somewhere accessible to the processing
layer. The dominant storage abstraction for this purpose in the 2010s
was the \term{data lake}: a centralised repository that stores raw data
at any scale and in any format---structured, semi-structured, and
unstructured---without requiring upfront schema definition.

The term was coined by James Dixon, then CTO of Pentaho, in 2010, in
contrast to the \term{data mart}: ``If you think of a data mart as a
store of bottled water---cleansed and packaged and structured for easy
consumption---the data lake is a large body of water in a more natural
state.''

%--------------------------------------------------------------------
\subsection{Data Lake Architecture}
\label{subsec:lake-architecture}
%--------------------------------------------------------------------

A well-designed data lake has four layers that transform raw ingested
bytes into query-ready, governed datasets.

\begin{defbox}{Data Lake Four-Layer Architecture}
\begin{enumerate}
  \item \textbf{Ingestion Layer.} Pulls data from all sources and
        writes raw files to the storage layer without transformation.
        Tools: Apache Sqoop (RDBMS bulk), Apache Flume (log streaming),
        Apache Kafka (event streaming), Apache NiFi (visual pipelines),
        REST API connectors.
  \item \textbf{Storage Layer.} Holds all data in its raw and processed
        forms. Organised into \emph{zones}:
        \begin{itemize}
          \item \textbf{Raw zone}: byte-for-byte copy of source data;
                never modified after ingest.
          \item \textbf{Processed zone}: cleaned, de-duplicated,
                schema-normalised data; typically Parquet or ORC.
          \item \textbf{Curated zone}: aggregated, business-ready
                datasets; typically Hive tables or Delta Lake tables.
        \end{itemize}
        Underlying storage: HDFS (on-premise) or S3/Azure ADLS
        (cloud-native).
  \item \textbf{Processing Layer.} Transforms and queries data on demand.
        Tools: Apache Spark (batch ETL and ML), Apache Hive (SQL
        analytics), Apache Presto/Trino (interactive queries), Apache Flink
        (stream processing).
  \item \textbf{Catalogue and Governance Layer.} Tracks schemas, data
        lineage, ownership, and access policies across all datasets.
        Tools: Apache Atlas, AWS Glue Data Catalog, Hive Metastore,
        Apache Ranger (access control).
\end{enumerate}
\end{defbox}

\begin{figure}[htbp]
\centering
\begin{tikzpicture}[
    box/.style={rectangle, rounded corners=4pt, minimum width=2.5cm,
                minimum height=0.75cm, font=\small\sffamily, align=center},
    layer/.style={rectangle, rounded corners=5pt, minimum width=10.0cm,
                  minimum height=1.2cm, font=\small\bfseries\sffamily,
                  align=center},
    arr/.style={-{Stealth[length=4pt,width=4pt]}, thick}
  ]

  %% Sources
  \node[box, fill=DEGray!15] (rdb)  at (-4.8, 7.0) {RDBMS};
  \node[box, fill=DEGray!15] (logs) at (-2.8, 7.0) {Logs};
  \node[box, fill=DEGray!15] (api)  at (-0.8, 7.0) {APIs};
  \node[box, fill=DEGray!15] (iot)  at (1.2, 7.0)  {IoT};
  \node[box, fill=DEGray!15] (files) at (3.2, 7.0) {Files};

  %% Ingestion layer
  \node[layer, fill=DEBlue!15, draw=DEBlue!35, minimum width=10.0cm]
        (ingest) at (-0.8, 5.8)
        {Ingestion Layer\\{\small Sqoop · Flume · Kafka · NiFi}};

  \foreach \src in {rdb,logs,api,iot,files} {
    \draw[arr, DEGray!50] (\src.south)--(\src.south |- ingest.north);
  }

  %% Storage layer: three zones
  \node[layer, fill=DERed!10, draw=DERed!25, minimum width=3.0cm,
        minimum height=1.8cm]
        (raw) at (-3.3, 4.2) {Raw Zone\\{\small byte-for-byte copy}};
  \node[layer, fill=DEAccent!12, draw=DEAccent!30, minimum width=3.0cm,
        minimum height=1.8cm]
        (proc) at (-0.8, 4.2) {Processed Zone\\{\small Parquet/ORC}};
  \node[layer, fill=DEGreen!12, draw=DEGreen!30, minimum width=3.0cm,
        minimum height=1.8cm]
        (cur) at (1.7, 4.2) {Curated Zone\\{\small Hive tables}};

  \draw[arr, DEBlue!50] (ingest.south)--(-3.3,5.3)--(-3.3,5.1);
  \draw[arr, DERed!40] (raw)--(proc);
  \draw[arr, DEAccent!50] (proc)--(cur);

  %% Processing layer
  \node[layer, fill=DEBlue!10, draw=DEBlue!25, minimum width=8.0cm]
        (processing) at (-0.8, 2.6)
        {Processing Layer\\{\small Spark · Hive · Presto · Flink}};

  \draw[arr, DEGreen!50] (cur.south)--(cur.south |- processing.north);
  \draw[arr, DEAccent!40] (proc.south)--(proc.south |- processing.north);

  %% Governance overlay
  \node[layer, fill=DEPurple!10, draw=DEPurple!30, minimum width=10.0cm]
        (gov) at (-0.8, 1.2)
        {Catalogue \& Governance Layer (Atlas · Ranger · Glue)};
  \draw[arr, DEPurple!40] (processing.south)--(gov.north);

\end{tikzpicture}
\caption{Data lake four-layer architecture. Raw data is ingested without
         schema enforcement, then progressively transformed through raw,
         processed, and curated zones. The governance layer enforces data
         quality, lineage tracking, and access control across all zones.}
\label{fig:data-lake}
\end{figure}

%--------------------------------------------------------------------
\subsection{The Data Swamp Problem and Governance}
\label{subsec:data-swamp}
%--------------------------------------------------------------------

A data lake without active governance degrades rapidly into a
\term{data swamp}: a repository where data accumulates without metadata,
owners, or quality guarantees, rendering it unusable in practice despite
being physically present.

Data swamp symptoms are predictable:
\begin{itemize}
  \item \textbf{Undiscoverability.} Analysts cannot find the dataset
        they need because there is no searchable catalogue. Directories
        like \texttt{/data/raw/2019/export\_v3\_FINAL\_FINAL2} proliferate
        without documentation.
  \item \textbf{Unknown freshness.} Without a timestamp or freshness
        SLA in the catalogue, an analyst cannot tell whether
        \texttt{/data/curated/customers} was updated today or six months ago.
  \item \textbf{Schema drift.} Multiple copies of the same dataset
        accumulated over time with inconsistent column names and types,
        because different teams ran their own transformations without
        a shared schema registry.
  \item \textbf{Privacy leakage.} Personally identifiable information
        (PII) mixes with non-sensitive data because ingest bypassed
        the governance layer. Under Bangladesh's Digital Security Act and
        Bangladesh Bank regulations on data localisation, this can result
        in regulatory penalties.
  \item \textbf{Unreliable quality.} Without automated quality checks
        at ingest time, downstream ML models silently train on corrupted
        or incomplete data.
\end{itemize}

\begin{systemdesign}{Data Lake Governance in Practice}
Preventing a data swamp requires enforcing governance at \emph{ingest
time}, not as a remediation activity. Four controls are foundational:
\begin{enumerate}
  \item \textbf{Mandatory catalogue registration.} Every ingest pipeline
        must register its output dataset in the data catalogue (Apache
        Atlas or AWS Glue) with at minimum: owner, source system,
        update frequency, schema version, and PII classification.
        Ingest pipelines that fail to register are blocked by a catalogue
        gateway.
  \item \textbf{Schema enforcement at the boundary.} Use Apache Avro with
        a Schema Registry for Kafka-sourced data, and enforce schema
        compatibility checks before writing to the processed zone. A
        schema-breaking change triggers a Slack alert to the owning team.
  \item \textbf{Automated data quality checks.} Deploy \tool{Apache
        Griffin} or \tool{Great Expectations} to run null-rate, outlier,
        and referential integrity checks on every batch landed in the
        processed zone. Datasets failing checks are quarantined in a
        \texttt{/data/quarantine/} zone and not promoted to curated.
  \item \textbf{Column-level access control.} Use Apache Ranger to define
        column-level access policies: analysts see masked customer names
        and phone numbers; only the compliance team sees raw PII.
        Row-level security restricts access to regional data.
\end{enumerate}
\end{systemdesign}

%--------------------------------------------------------------------
\subsection{Table Formats: Delta Lake and Apache Iceberg}
\label{subsec:table-formats}
%--------------------------------------------------------------------

Object stores (HDFS, S3) do not natively support ACID transactions:
two concurrent Spark jobs writing to the same directory can corrupt
each other's output (the \emph{lost-update} problem). Traditional
partitioned Hive tables suffer from further limitations: schema
evolution requires \texttt{ALTER TABLE} operations that can break
existing readers; time-travel (querying the state of a table at a
past timestamp) is not supported; and partition discovery is an
expensive NameNode scan for tables with thousands of partitions.

\term{Open table formats}---specifically \tool{Delta Lake} (Databricks,
2019) and \tool{Apache Iceberg} (Netflix, 2020)---solve these
problems by adding a \emph{transaction log} layer on top of Parquet
files in an object store.

\begin{defbox}{Transaction Log in Open Table Formats}
An \textbf{open table format} such as Delta Lake or Apache Iceberg
maintains a \emph{transaction log}: an ordered list of committed
operations (add file, remove file, update schema, change partition
spec). Each operation is an atomic JSON or Avro record written to a
separate log directory. Properties enabled by the transaction log:
\begin{itemize}
  \item \textbf{ACID writes.} A Spark job writes new Parquet files to
        a staging location, then atomically commits a ``add files''
        entry to the transaction log. If the job fails, no partial
        state is visible to readers.
  \item \textbf{Snapshot isolation.} Readers always see a consistent
        snapshot of the table as of the latest committed transaction;
        concurrent writers do not interfere.
  \item \textbf{Time travel.} Querying the table \texttt{AS OF TIMESTAMP
        '2024-01-01'} reads the transaction log backward to find the
        snapshot valid at that time, then reads the corresponding
        Parquet files.
  \item \textbf{Schema evolution.} New columns, renamed columns, and
        changed nullability are recorded in the log and propagated to
        readers automatically.
\end{itemize}
\end{defbox}

\begin{lstlisting}[style=pyspark, caption={Delta Lake: ACID writes, time travel, and schema evolution}]
from delta.tables import DeltaTable
from pyspark.sql import SparkSession
from pyspark.sql.functions import col, current_timestamp

spark = (SparkSession.builder
         .appName("DeltaLakeDemo")
         .config("spark.sql.extensions",
                 "io.delta.sql.DeltaSparkSessionExtension")
         .config("spark.sql.catalog.spark_catalog",
                 "org.apache.spark.sql.delta.catalog.DeltaCatalog")
         .getOrCreate())

TABLE_PATH = "s3://data-lake/curated/transactions"

# --- Write (ACID-guaranteed) ---
df = spark.read.parquet("s3://data-lake/raw/transactions/2024-06-28/")
(df.write
   .format("delta")
   .mode("append")
   .partitionBy("txn_date", "region")
   .save(TABLE_PATH))

# --- Merge (upsert: update existing rows, insert new ones) ---
target = DeltaTable.forPath(spark, TABLE_PATH)
updates = spark.read.parquet("s3://data-lake/raw/corrections/")

(target.alias("t")
       .merge(updates.alias("u"),
              "t.txn_id = u.txn_id")
       .whenMatchedUpdateAll()
       .whenNotMatchedInsertAll()
       .execute())

# --- Time travel: read the table as it was yesterday ---
yesterday = spark.read.format("delta") \
                 .option("timestampAsOf", "2024-06-27") \
                 .load(TABLE_PATH)
yesterday.count()

# --- Inspect transaction history ---
(DeltaTable.forPath(spark, TABLE_PATH)
           .history()
           .select("version", "timestamp", "operation", "operationMetrics")
           .show(10, truncate=False))
\end{lstlisting}

\begin{keyinsight}{Delta Lake and Iceberg Converge the Lake and Warehouse}
The combination of (1) object-store economics (\$0.02/GB/month on S3 vs.\
\$2--\$5/GB/month for proprietary data warehouse storage), (2) ACID
transaction semantics from Delta Lake or Iceberg, and (3) SQL query
engines like Presto and Spark SQL creates a new architecture called the
\term{Lakehouse} (Armbrust et al., 2020)---a data lake with the ACID
guarantees and performance characteristics of a data warehouse, without
the data movement and duplication costs of maintaining separate systems.
The Lakehouse is the subject of Chapter~\ref{ch:modern-stack}.
\end{keyinsight}

%%====================================================================
\section{Research Spotlight}
\label{sec:ch8-research}
%%====================================================================

\begin{researchspotlight}{Kreps, Narkhede \& Rao (2011) --- Kafka: A
Distributed Messaging System for Log Aggregation}

\textbf{Full citation:} Kreps, J., Narkhede, N., \& Rao, J. (2011).
Kafka: A Distributed Messaging System for Log Aggregation. In
\textit{Proceedings of the NetDB Workshop at VLDB 2011}. Seattle, WA.
Cited by $>$3{,}500 (Google Scholar, 2024).

\medskip
\textbf{Problem context.} By 2010, LinkedIn was generating billions of
activity events per day (profile views, connection requests, job searches,
feed impressions). The existing pipeline---custom ad-hoc processes copying
log files over NFS to Hadoop---was fragile, could not replay past data,
and could not scale to the growing event volume. Existing message brokers
(ActiveMQ, RabbitMQ) were designed for transactional messaging at thousands
of messages/sec, not streaming at 100~MB/s per broker with multi-day retention.

\medskip
\textbf{Key insight.} \textit{``A log is the most primitive form of storage.
If you model your data infrastructure as a distributed, replicated,
append-only log, you get durability, replayability, and decoupling for free.''}
This insight---that the database write-ahead log primitive could be
generalised into a distributed messaging infrastructure---is the paper's
central contribution.

\medskip
\textbf{Technical contributions.}
\begin{enumerate}
  \item \textbf{Topic-partition-offset data model.} The paper formalises
        the abstraction of a topic (named log), partition (ordered,
        independent sub-log for parallelism), and offset (consumer-tracked
        position). This model provides ordered delivery within a partition,
        unlimited consumer fan-out (each consumer group tracks its own
        offsets independently), and log-time ordering without a global
        coordinator.
  \item \textbf{Sequential disk I/O and zero-copy.} The paper quantifies
        the performance advantage of append-only log files over B-tree
        indexed stores: sequential disk bandwidth on spinning media is
        600~MB/s versus 100~MB/s for random access. Kafka's \texttt{sendfile()}
        zero-copy transfer reduces consumer-side CPU to near zero at peak
        throughput. The benchmark reports 505{,}000 producer msgs/sec and
        220{,}000 consumer msgs/sec on a single broker, outperforming
        ActiveMQ by $11.7\times$ and $5.5\times$ respectively.
  \item \textbf{Pull-based consumers.} Unlike push-based brokers that must
        track per-consumer queue state, Kafka's brokers maintain only the
        log itself. Consumers pull messages at their own rate; the broker
        never needs to know which consumers exist or what they have consumed.
        This eliminates a significant bottleneck: a slow consumer cannot
        slow down the broker.
\end{enumerate}

\medskip
\textbf{Impact.} Kafka was open-sourced to the Apache Software Foundation
in 2012. By 2021, LinkedIn's Kafka cluster processed 1.8~trillion
messages/day across 4{,}400 brokers; production deployments at Netflix,
Uber, Airbnb, Twitter/X, and thousands of enterprises handle aggregate
event volumes exceeding 10~trillion messages/day. The paper triggered
a rethinking of data infrastructure: Jay Kreps's 2013 essay \textit{The
Log} (which extends the paper's ideas) became one of the most widely
shared technical articles in the data engineering community. The paper
directly inspired the Dataflow model (Akidau et al., 2015) and the
Kappa architecture (Kreps, 2014), both covered in Chapter~\ref{ch:modern-stack}.
\end{researchspotlight}

%%====================================================================
\section{Case Study: Netflix Keystone Real-Time Data Pipeline}
\label{sec:netflix-case}
%%====================================================================

\begin{casestudy}{Netflix Keystone Pipeline --- 100 Billion Events per Day}

\textbf{Context and scale.} Netflix serves 200 million subscribers
(as of 2021) across 190 countries. Every user action---play, pause,
seek, quality switch, UI click, search query, content impression---is
recorded as an event. At 1.3~million events/second at peak (Friday at
9\,pm US Eastern) and 100~billion events/day at average load, Netflix's
data infrastructure operates at a scale few organisations approach.
The Keystone pipeline is the internal name for the real-time event
ingestion and routing system that processes all of these events before
they reach analytics, personalisation, A/B testing, and operational
monitoring systems.

\medskip
\textbf{Problem.} Before Keystone, Netflix used a custom log-based
pipeline where applications wrote events to local disk, a daemon
collected them and pushed to a queue, and batch jobs processed the
queue nightly. This pipeline could not support real-time personalisation
(the home screen must reflect what a user watched 10 minutes ago, not
yesterday) or same-day A/B test results.

\medskip
\textbf{Architecture.}

Netflix's Keystone architecture is a hybrid Lambda-Kappa system:

\begin{center}
\small
\renewcommand{\arraystretch}{1.3}
\begin{tabular}{llll}
  \toprule
  \textbf{Layer} & \textbf{Tool} & \textbf{Latency} & \textbf{Output} \\
  \midrule
  Event bus
    & Apache Kafka (4{,}000+ topics)
    & $<$100\,ms
    & All raw events \\
  Speed layer
    & Apache Flink (150+ jobs)
    & $<$1\,min
    & Real-time dashboards, alerts \\
  Batch layer
    & Kafka $\to$ S3 $\to$ Spark (EMR) + Iceberg
    & $\sim$1\,hour--1\,day
    & ML datasets, billing, BI \\
  Serving
    & Cassandra (user state), Druid (OLAP)
    & $<$10\,ms
    & Personalisation, analytics \\
  \bottomrule
\end{tabular}
\end{center}

\medskip
Every device (TV, mobile, web browser) sends events to Kafka topics via
a client SDK. Kafka acts as the \emph{single event bus}: no downstream
system receives events directly from devices. Flink jobs subscribe to
Kafka topics and compute real-time aggregates (per-user play counts,
per-title impression rates, A/B test counters). A Kafka Connect S3
Sink connector writes all raw events to Amazon S3 in Avro format,
partitioned by hour. Spark jobs on Amazon EMR read from S3, compute
higher-fidelity aggregates using Apache Iceberg ACID tables, and
write results to the curated S3 zone.

\medskip
\textbf{A/B testing pipeline.} Netflix runs $>$200 concurrent A/B tests
at any time---different thumbnail artwork, autoplay delays, recommendation
algorithm variants, and UI layouts. Each test requires:
\begin{enumerate}
  \item An \emph{assignment event} when the user is bucketed into control
        or treatment (stored in Cassandra keyed by user ID).
  \item An \emph{exposure event} logged to Kafka when the user actually
        sees the tested feature.
  \item \emph{Outcome events} (play, click, subscription) logged to
        Kafka within seconds.
\end{enumerate}
A Flink job performs a three-way event-time window join over the
\texttt{assignments}, \texttt{exposures}, and \texttt{outcomes} Kafka
topics, computing conversion rates with a 30-second watermark lag to
handle late-arriving mobile events. Experiment dashboards are updated
every 5 minutes, enabling engineers to ship winning variants the same day
instead of waiting for a nightly batch run.

\medskip
\textbf{Engineering lessons.}
\begin{enumerate}
  \item \textbf{Kafka as the single source of truth eliminates
        inconsistency.} By having both the speed layer (Flink) and the
        batch layer (Spark on S3) read from the same Kafka topics, Netflix
        avoids the classic Lambda problem of two separate ingestion paths
        diverging. The S3 landing is just a materialised snapshot of the
        Kafka log.
  \item \textbf{7-day retention enables operational recovery.} When a
        Flink job has a bug, engineers reset the consumer group offset to
        the time before the bug was introduced and reprocess---without
        contacting source teams or rebuilding input datasets.
  \item \textbf{Schema Registry is non-negotiable at 4{,}000 topics.}
        Netflix co-developed the Confluent Schema Registry and mandates
        Avro schema registration for every Kafka topic. A schema change
        that would break downstream consumers is rejected at deployment time.
  \item \textbf{Backpressure is the production problem, not throughput.}
        During Super Bowl commercial breaks, event rates spike $5\times$
        in 30 seconds. Kafka absorbs the burst; Flink consumers catch
        up within minutes because they are stateless with respect to the
        Kafka log. Designing for backpressure---not peak throughput---is
        the key operational insight.
  \item \textbf{Iceberg enables schema evolution without downtime.}
        Netflix's batch pipeline manages 5{,}000+ Iceberg tables, many
        of which evolve weekly. Iceberg's schema evolution support allows
        columns to be added or renamed without rewriting existing Parquet
        files or pausing consuming Spark jobs.
\end{enumerate}

\medskip
\textbf{Open-source contributions.} Netflix open-sourced \tool{Metacat}
(a unified metadata store and catalogue for Hive, Iceberg, S3, and Druid),
and was a co-developer of \tool{Apache Iceberg}, both driven by the
requirements of the Keystone pipeline. Netflix's engineering blog posts
on Kafka, Flink, and Iceberg (2015--2022) are widely used reference
material in the data engineering community.
\end{casestudy}

%%====================================================================
\section*{Chapter Summary}
\label{sec:ch8-summary}
%%====================================================================

This chapter traced the evolution of data integration from fragile
point-to-point ETL to the continuous, event-driven pipelines that
underpin modern data infrastructure.

The $N^2$ coupling problem of direct ETL connections is resolved by a
central \emph{event bus}, of which Apache Kafka is the dominant
implementation. Kafka's performance advantage---$10\times$ throughput over
traditional message brokers---derives from two engineering decisions:
sequential-write-only disk access and zero-copy transfer via \texttt{sendfile()}.
Its durability advantage---consumers can replay any historical window---derives
from the append-only log abstraction that treats messages as records in a
database, not as transient queue entries.

Two architectural patterns organise batch and stream processing around
Kafka: the \emph{Lambda architecture} (Marz, 2011), which serves both
accurate batch views and low-latency real-time views at the cost of
maintaining two code paths; and the \emph{Kappa architecture} (Kreps, 2014),
which eliminates the batch layer by using Kafka log replay for historical
reprocessing, at the cost of higher storage requirements and the constraint
that the stream processing engine must be capable of reprocessing history
at acceptable speed.

Data lakes provide the petabyte-scale storage substrate for pipeline output.
Without active governance---schema registration, data quality checks,
ownership tracking, and access control---a data lake degrades into a
data swamp. Open table formats (Delta Lake, Apache Iceberg) add ACID
transaction semantics to object-store data lakes, enabling time travel,
schema evolution, and concurrent writer safety. They form the foundation
of the Lakehouse pattern explored in Chapter~\ref{ch:modern-stack}.

%%====================================================================
\section*{Exercises}
\label{sec:ch8-exercises}
%%====================================================================

\exercisesection{Short Questions}

\sq What is the $N^2$ problem in point-to-point ETL integration? With 30
systems, how many direct connections are required in the worst case? How
does an event bus reduce this number?

\sq Describe the three components of an Apache Flume agent. What is the
difference between a Memory Channel and a File Channel in terms of
durability?

\sq What is the difference between an Apache Kafka \emph{topic} and a
Kafka \emph{partition}? Why must messages with the same key be routed to
the same partition?

\sq Explain what a Kafka \emph{offset} is. Who is responsible for tracking
offsets in the Kafka data model---the broker or the consumer? What is the
benefit of this design?

\sq What is a Kafka consumer group? If a topic has 6 partitions and a
consumer group has 4 consumers, how many partitions does each consumer
process (assuming uniform assignment)?

\sq Explain the Linux \texttt{sendfile()} system call and its role in
Kafka's consumer throughput. How many memory copies does a traditional
message broker require compared to Kafka?

\sq Describe the Lambda architecture's three layers. What is the purpose
of the Serving Layer?

\sq What is the ``dual-code-path problem'' in the Lambda architecture?
Why does it create a maintenance burden?

\sq In the Kappa architecture, how is historical reprocessing accomplished?
Describe the step-by-step reprocessing workflow when business logic
changes.

\sq What is a data lake? List four symptoms that indicate a data lake has
become a data swamp.

\sq What is an open table format (e.g., Delta Lake or Apache Iceberg)?
Name three features that open table formats add to a standard Parquet-on-S3
data lake.

\sq Compare Apache Sqoop and Apache Kafka on: (a) data direction, (b) latency,
(c) replay capability, and (d) the underlying parallelism mechanism.

\exercisesection{Analytical Problems}

\ap \textbf{Kafka Partition Sizing.}
A mobile payment company (similar to bKash) expects the following workload
on its \texttt{wallet-transactions} Kafka topic:
\begin{itemize}
  \item Peak producer throughput: 80{,}000 transactions/second
  \item Average message size: 512 bytes
  \item Required consumer throughput per partition: 10~MB/s (broker limit)
  \item Replication factor: 3
  \item Required retention: 14 days
  \item Average message rate: 20{,}000 transactions/second
\end{itemize}
(a) Calculate the minimum number of partitions required to sustain the
peak producer throughput. (b) Calculate the total storage required on the
Kafka cluster for 14-day retention (assuming average rate, before replication).
(c) Calculate the total storage \emph{after} replication. (d) If each
broker node has 12~TB of usable disk, how many brokers are needed?

\ap \textbf{Lambda vs.\ Kappa Architecture Selection.}
A Bangladeshi commercial bank is designing a data pipeline for three
workloads:
\begin{enumerate}
  \item Real-time fraud detection: flag suspicious card-swipe transactions
        within 200~ms with a rolling 24-hour transaction history per card.
  \item Monthly regulatory capital report: aggregates 36 months of
        transaction data into 47 statutory metrics as required by
        Bangladesh Bank Circular BRPD~18.
  \item Live risk dashboard: shows current total exposure per counterparty
        updated every 30 seconds, combining live transactions (from the
        past 5 minutes) with a settled position snapshot from the previous
        close of business.
\end{enumerate}
For each workload, identify whether Lambda, Kappa, or a different
architecture is most appropriate. Justify each choice with reference
to latency requirement, reprocessing cost, and regulatory constraints.

\ap \textbf{Consumer Group Rebalancing.}
A Kafka topic has 12 partitions. A consumer group starts with 6 consumers,
each processing 2 partitions.
\begin{enumerate}[(a)]
  \item Consumer~3 fails (crashes without a graceful shutdown). Describe
        what Kafka does: which consumers take over its partitions, and
        how long does the rebalance take (Kafka's default
        \texttt{session.timeout.ms} is 45 seconds)?
  \item Two new consumers join the group simultaneously. How are the 12
        partitions reassigned under the default RangeAssignor strategy?
  \item A new version of the consumer is deployed with 12 instances
        (one per partition). What happens if 4 additional instances are
        accidentally started, making the total 16?
\end{enumerate}

\ap \textbf{Data Lake Zone Design.}
Design the zone structure for a national healthcare data lake that will
store:
\begin{itemize}
  \item Patient records from 64 district hospitals (structured, RDBMS)
  \item Doctor appointment audio recordings (unstructured, MP3)
  \item Lab result CSV uploads from 200 diagnostic centres
  \item Real-time vital signs from 1{,}000 ICU monitoring devices (Kafka)
\end{itemize}
For each data source: (a) identify which ingestion tool to use (Sqoop,
Flume, Kafka, NiFi); (b) specify the landing zone (raw, processed, or
curated); (c) identify the format in the processed zone; (d) describe
which governance controls are mandatory given that patient records are
PII under the National Health Policy 2011 (Bangladesh).

\exercisesection{Design Problems}

\dpr \textbf{bKash Real-Time Fraud Pipeline.}
Design a complete Kafka-based data pipeline for bKash (Bangladesh's
largest mobile financial service, 50 million users, 10 million
transactions/day) that must:
\begin{enumerate}[(a)]
  \item Detect fraudulent transactions in real time ($<$200~ms decision
        latency) using velocity checks (more than 5 transactions from
        the same wallet in 60 seconds).
  \item Power a live operations dashboard showing total transaction
        volume and flagged transactions per minute, updated every 30
        seconds.
  \item Provide a daily reconciliation file to Bangladesh Bank containing
        all flagged transactions and final dispositions (confirmed fraud
        vs.\ false positive).
  \item Allow the fraud detection rules to be updated and back-tested
        against 30 days of historical transactions without re-ingesting
        from the source database.
\end{enumerate}
Your design must specify: Kafka topic structure (topic names, partition
count, key, retention); the processing framework for each workload (Kafka
Streams, Flink, Spark batch); which architecture pattern (Lambda or Kappa)
fits each requirement and why; and the serving layer technology for the
dashboard and the reconciliation file.

\dpr \textbf{University Data Lake Design.}
Premier University, Chittagong, wishes to build a data lake to unify
its operational data for academic analytics, student counselling, and
research reporting. Data sources include: student RDBMS (Oracle, 50K
student records), LMS (Moodle, JSON event logs), library system (MySQL,
book borrowing records), exam results (Excel uploads by faculty), and
Wi-Fi access logs (syslog, 100K events/day).
\begin{enumerate}[(a)]
  \item Design the ingestion pipeline for each source.
  \item Define the zone structure and the schema for the curated zone
        \texttt{student\_engagement} table (columns, types, partitioning).
  \item Describe three data quality checks that must pass before a dataset
        is promoted from raw to processed zone.
  \item Identify which datasets contain PII under the Bangladesh
        Digital Security Act and describe the access control policies
        (who can see what, in what form) to be enforced by Apache Ranger.
\end{enumerate}

\exercisesection{Programming Exercises}

\pe \textbf{Kafka Producer and Consumer (Python).}
Using the \texttt{kafka-python} library and a local single-broker
Docker container:
\begin{enumerate}[(a)]
  \item Implement a producer that generates 50{,}000 synthetic mobile
        payment events (fields: \texttt{txn\_id}, \texttt{sender},
        \texttt{receiver}, \texttt{amount\_bdt}, \texttt{timestamp\_ms}).
        Key each message by \texttt{sender} to ensure all transactions
        from a given sender land in the same partition.
  \item Implement a consumer in group \texttt{analytics-v1} that reads
        all messages and computes: (i) total transaction volume in BDT;
        (ii) top-10 senders by count; (iii) average transaction amount
        per minute over the message timestamps.
  \item Reset the consumer offset to the beginning and re-run the
        computation. Verify that results are identical (idempotency test).
  \item Create a second consumer group \texttt{fraud-v1}. Verify that
        both groups independently read all messages from offset~0,
        with no offset interference between them.
\end{enumerate}

\begin{lstlisting}[style=pyspark, caption={PE8.1: Kafka producer skeleton}]
import json, random, time, uuid
from kafka import KafkaProducer

BROKER = "localhost:9092"
TOPIC  = "mobile-payments"

producer = KafkaProducer(
    bootstrap_servers=BROKER,
    key_serializer=str.encode,
    value_serializer=lambda v: json.dumps(v).encode(),
    acks='all',
    retries=3,
)

SENDERS = [f"01{random.randint(3,9)}{random.randint(10000000,99999999)}"
           for _ in range(1000)]

for _ in range(50_000):
    sender = random.choice(SENDERS)
    event = {
        "txn_id":       str(uuid.uuid4()),
        "sender":       sender,
        "receiver":     random.choice(SENDERS),
        "amount_bdt":   round(random.uniform(10, 10_000), 2),
        "timestamp_ms": int(time.time() * 1000),
    }
    producer.send(TOPIC, key=sender, value=event)

producer.flush()
print("Done: 50,000 events produced")
\end{lstlisting}

\pe \textbf{Lambda Architecture Simulator (PySpark).}
Simulate a Lambda architecture for a ride-sharing app (similar to Pathao)
using PySpark and local files:
\begin{enumerate}[(a)]
  \item \textbf{Batch layer:} Read a 1-million-row Parquet file of
        historical ride events (columns: \texttt{ride\_id},
        \texttt{driver\_id}, \texttt{pickup\_zone}, \texttt{fare\_bdt},
        \texttt{ride\_date}). Compute the hourly fare total per pickup
        zone and write as a batch view Parquet file.
  \item \textbf{Speed layer:} Generate 10{,}000 synthetic ``live'' events
        for the current hour (using \texttt{datetime.now()}). Compute
        the same fare-total-per-zone aggregate as an in-memory DataFrame.
  \item \textbf{Serving layer:} Union the batch view with the real-time
        view on \texttt{(pickup\_zone, hour)} and resolve overlaps (the
        current hour appears in both; the batch view value should be
        discarded in favour of the real-time value).
  \item Report the total number of rides and fare sum per zone for the
        most recent 24-hour period.
\end{enumerate}

\begin{lstlisting}[style=pyspark, caption={PE8.2: Lambda serving layer union in PySpark}]
from pyspark.sql import SparkSession
from pyspark.sql.functions import col, sum as _sum, hour, lit
from datetime import datetime, timedelta

spark = SparkSession.builder.appName("LambdaSim").getOrCreate()

# --- Batch layer view (pre-computed) ---
batch_view = spark.read.parquet("data/batch_views/fare_by_zone_hour/")
# Schema: pickup_zone, hour_ts, total_fare_bdt, source='batch'

# --- Speed layer (in-memory simulation) ---
import random

live_events = spark.createDataFrame([
    {
        "ride_id":      f"R{i:06d}",
        "driver_id":    random.randint(1, 5000),
        "pickup_zone":  random.choice(["Gulshan","Dhanmondi","Mirpur",
                                       "Uttara","Motijheel","Wari"]),
        "fare_bdt":     round(random.uniform(50, 800), 2),
        "ride_ts":      datetime.now() - timedelta(seconds=random.randint(0, 3599)),
    }
    for i in range(10_000)
])

real_time_view = (live_events
    .groupBy("pickup_zone")
    .agg(_sum("fare_bdt").alias("total_fare_bdt"))
    .withColumn("hour_ts", lit(datetime.now().replace(minute=0, second=0,
                                                       microsecond=0)))
    .withColumn("source", lit("realtime")))

# --- Serving layer: union, prefer real-time for current hour ---
current_hour = datetime.now().replace(minute=0, second=0, microsecond=0)
batch_historical = batch_view.filter(col("hour_ts") < lit(current_hour))

served = batch_historical.unionByName(real_time_view)
served.orderBy("pickup_zone", "hour_ts").show(20, truncate=False)
\end{lstlisting}

\pe \textbf{Data Quality Gate (Python + Great Expectations).}
Implement an automated data quality gate that must pass before a dataset
is promoted from the raw zone to the processed zone of a data lake:
\begin{enumerate}[(a)]
  \item Load a CSV file of customer records (fields: \texttt{customer\_id},
        \texttt{mobile}, \texttt{name}, \texttt{district}, \texttt{balance\_bdt},
        \texttt{account\_opened}).
  \item Using the \texttt{great\_expectations} library, define and run
        the following expectations:
        \begin{itemize}
          \item \texttt{customer\_id} must not be null and must be unique.
          \item \texttt{mobile} must match the regex
                \texttt{\textasciicircum{}01[3-9]\textbackslash d\{8\}\$}
                (Bangladeshi mobile number format).
          \item \texttt{balance\_bdt} must be between 0 and 10{,}000{,}000.
          \item \texttt{district} must be one of the 64 districts of Bangladesh.
          \item No more than 1\% of rows may have a null \texttt{name}.
        \end{itemize}
  \item If all expectations pass, write the dataset to the processed zone
        as Parquet. If any expectation fails, write the dataset to a
        quarantine directory and raise an alert (print a summary report).
\end{enumerate}
%%====================================================================
\section*{Further Reading}
\label{sec:ch8-further}
%%====================================================================

\begin{itemize}
  \item \textbf{\cite{Kafka2011}} --- The original Kafka paper. Short and
        readable; the benchmark section is directly relevant to
        understanding why Kafka displaced existing message brokers.
  \item \textbf{\cite{Kleppmann2017}} --- Chapters 10 and 11 cover batch
        processing (Lambda batch layer) and stream processing (speed layer
        and Kappa) at graduate level. The ``unified log'' discussion
        expands Kreps's original insight into a full architectural theory.
  \item \textbf{\cite{White2015}} --- Chapters 14 (Flume) and 15 (Sqoop)
        provide implementation details not covered in this chapter, including
        Flume interceptors, channel selectors, and Sqoop's \texttt{sqoop-import-all-tables}
        command.
  \item \textbf{Kreps, J. (2013). ``The Log: What Every Software Engineer
        Should Know About Real-Time Data's Unifying Abstraction.''
        LinkedIn Engineering Blog.} Freely available online; the most
        important essay in data engineering of the last decade.
  \item \textbf{Arun, M.\ et al.\ (2016). ``Kafka Inside Keystone Pipeline.''
        Netflix Tech Blog.} The primary Netflix source describing the
        Keystone architecture, Kafka configuration, and engineering
        challenges at 1.3~million events/second.
\end{itemize}
